% =====================================================================
%  PROTOCOLO EXPERIMENTAL -- Proyecto SIMPA
%  Enfoque 1: calidad de requisitos elicitados por analistas humanos
%             frente a los generados por un modelo grande de lenguaje
%  Equipo AHMRV -- ISR-401 -- UTEQ -- 2026-2027 PPA
%  Compilacion: pdflatex protocolo.tex   (dos veces)
% =====================================================================
\documentclass[11pt,a4paper]{article}
\usepackage[utf8]{inputenc}
\usepackage[T1]{fontenc}
\usepackage{amsmath}
\usepackage[margin=2.4cm]{geometry}
\usepackage{longtable}
\usepackage{array}
\usepackage{xcolor}
\usepackage{colortbl}
\usepackage{enumitem}
\usepackage{fancyhdr}
\usepackage{lastpage}
\usepackage{titlesec}
\usepackage{url}
\usepackage[hidelinks]{hyperref}
\usepackage{caption}

\definecolor{uteqverde}{RGB}{0,102,51}
\definecolor{grisclaro}{RGB}{242,242,242}
\newcolumntype{L}[1]{>{\raggedright\arraybackslash}p{#1}}
\newcolumntype{C}[1]{>{\centering\arraybackslash}p{#1}}
\newcommand{\id}[1]{\texttt{#1}}
\renewcommand{\tablename}{Tabla}
\sloppy

\pagestyle{fancy}
\fancyhf{}
\lhead{\small Protocolo experimental --- SIMPA}
\rhead{\small Equipo AHMRV --- ISR-401}
\cfoot{\small Página \thepage\ de \pageref{LastPage}}
\renewcommand{\headrulewidth}{0.4pt}
\titleformat{\section}{\normalfont\large\bfseries\color{uteqverde}}{\thesection.}{0.6em}{}

\begin{document}

\begin{center}
{\large\bfseries UNIVERSIDAD TÉCNICA ESTATAL DE QUEVEDO\par}
{\normalsize Facultad de Ciencias de la Computación --- Carrera de Software\par}
\vspace{0.5cm}
\rule{\textwidth}{1pt}\par
\vspace{0.3cm}
{\LARGE\bfseries\color{uteqverde} PROTOCOLO EXPERIMENTAL\par}
\vspace{0.25cm}
{\large Calidad de los requisitos funcionales elicitados por analistas humanos\\
frente a los generados por un modelo grande de lenguaje\\
a partir del mismo material fuente\par}
\vspace{0.3cm}
\rule{\textwidth}{1pt}\par
\end{center}

\vspace{0.3cm}
\begin{longtable}{|L{4.0cm}|L{10.8cm}|}
\hline
\rowcolor{grisclaro}\multicolumn{2}{|l|}{\textbf{Identificación}}\\
\hline
\textbf{Proyecto} & SIMPA --- Sistema Inteligente de Mantenimiento de Palma Africana \\ \hline
\textbf{Equipo} & AHMRV --- Paralelo 4to ``A'' \\ \hline
\textbf{Enfoque} & Enfoque 1 de la Sección 5 de la guía de la Entrega 3 (2A) \\ \hline
\textbf{Diseño} & Cuasi-experimento con medidas apareadas y evaluación a ciegas \\ \hline
\textbf{Registro previo} & Marco de Ciencia Abierta (OSF) --- véase \id{osf\_registration.pdf} \\ \hline
\textbf{Estado} & Diseñado y registrado. \textbf{Ejecución pendiente} \\ \hline
\textbf{Versión} & 1.0 --- 3 de agosto de 2026 \\ \hline
\end{longtable}

% =====================================================================
\section{Motivación y contexto}

La revisión sistemática de Cheng \emph{et al.} sobre inteligencia artificial generativa aplicada a
la ingeniería de requisitos señala que la elicitación es una de las fases con menor evidencia
empírica acumulada. Las guías de Vogelsang y Fischbach para el uso de modelos grandes de lenguaje en
tareas de procesamiento de lenguaje natural aplicadas a requisitos proponen prácticas que aún no han
sido validadas en contextos latinoamericanos ni en dominios agroindustriales.

Este estudio aporta un caso de evidencia adicional: compara el conjunto de requisitos funcionales
elicitados por un equipo humano de estudiantes con el conjunto que produce un modelo grande de
lenguaje a partir de la misma transcripción de entrevista, evaluados a ciegas por personas expertas
independientes.

% =====================================================================
\section{Pregunta de investigación en formato PICOC}

\begin{longtable}{|C{2.6cm}|L{12.2cm}|}
\hline
\rowcolor{grisclaro}\textbf{Elemento} & \textbf{Definición para este estudio} \\
\hline
\endhead
\textbf{Población} & Requisitos funcionales de un sistema agroindustrial real de gestión y monitoreo de cultivo de palma africana \\ \hline
\textbf{Intervención} & Generación de requisitos funcionales mediante un modelo grande de lenguaje a partir de una transcripción de entrevista anonimizada \\ \hline
\textbf{Comparación} & Requisitos funcionales elicitados por el equipo humano a partir del mismo material fuente \\ \hline
\textbf{Resultado} & Puntuación en cinco dimensiones de calidad: completitud, ausencia de ambigüedad, verificabilidad, corrección respecto de la fuente y consistencia interna \\ \hline
\textbf{Contexto} & Proyecto académico de ingeniería de requisitos, Ecuador, con evidencia de campo propia \\ \hline
\caption{Marco PICOC del estudio}
\end{longtable}

\noindent
\textbf{RQ1.} ¿En qué dimensiones de calidad difieren los requisitos funcionales elicitados por un
equipo humano de los generados por un modelo grande de lenguaje a partir del mismo material fuente?

\noindent
\textbf{RQ2.} ¿Cuál es el grado de acuerdo entre personas evaluadoras independientes al valorar la
calidad de requisitos funcionales sin conocer su procedencia?

% =====================================================================
\section{Hipótesis}

Para cada dimensión de calidad $d \in \{$completitud, ausencia de ambigüedad, verificabilidad,
corrección, consistencia$\}$:

\begin{itemize}[nosep,leftmargin=1.6em]
  \item $H_{0,d}$: no existe diferencia entre la puntuación media del conjunto humano y la del
        conjunto generado por el modelo. Formalmente, $\mu_{H,d} = \mu_{L,d}$.
  \item $H_{1,d}$: existe diferencia estadísticamente significativa. Formalmente,
        $\mu_{H,d} \neq \mu_{L,d}$.
\end{itemize}

\noindent
La hipótesis alternativa se formula de forma \textbf{bilateral}. El equipo no anticipa la dirección
del efecto: la literatura reporta resultados mixtos, y postular una dirección sin base sería un
sesgo de expectativa.

% =====================================================================
\section{Variables}

\begin{longtable}{|C{2.6cm}|L{5.4cm}|L{6.8cm}|}
\hline
\rowcolor{grisclaro}\textbf{Tipo} & \textbf{Variable} & \textbf{Operacionalización} \\
\hline
\endhead
Independiente & Origen del requisito & Nominal dicotómica: humano / modelo de lenguaje \\ \hline
Dependiente & Completitud & Ordinal 1--5: presencia de los ocho atributos de la plantilla \\ \hline
Dependiente & Ausencia de ambigüedad & Ordinal 1--5: unicidad de interpretación del enunciado \\ \hline
Dependiente & Verificabilidad & Ordinal 1--5: existencia de un criterio objetivamente comprobable \\ \hline
Dependiente & Corrección respecto de la fuente & Ordinal 1--5: fidelidad al contenido de la transcripción \\ \hline
Dependiente & Consistencia interna & Ordinal 1--5: ausencia de contradicción con otros requisitos del mismo conjunto \\ \hline
Control & Material fuente & Idéntico para ambos conjuntos: transcripción de \id{ENTR-04} \\ \hline
Control & Rúbrica de evaluación & Idéntica para ambos conjuntos \\ \hline
Control & Ceguera & Las personas evaluadoras desconocen la procedencia de cada requisito \\ \hline
Control & Orden de presentación & Aleatorizado para evitar efecto de posición \\ \hline
\caption{Variables del estudio}
\end{longtable}

% =====================================================================
\section{Población, muestra y unidad de análisis}

\textbf{Unidad de análisis:} el requisito funcional individual.

\begin{longtable}{|L{5.6cm}|C{2.6cm}|L{6.0cm}|}
\hline
\rowcolor{grisclaro}\textbf{Conjunto} & \textbf{Tamaño} & \textbf{Procedencia} \\
\hline
\endhead
Conjunto A --- generado por el modelo & $\geq 25$ & Producido a partir de la transcripción anonimizada de \id{ENTR-04} \\ \hline
Conjunto B --- elicitado por el equipo & $\geq 25$ & Subconjunto del ERS trazado a \id{EV-04} y evidencias relacionadas \\ \hline
Personas evaluadoras & 3 (mínimo) & Estudiantes de otro paralelo de la misma asignatura, sin participación en el proyecto \\ \hline
\caption{Conjuntos y participantes}
\end{longtable}

\noindent
\textbf{Nota sobre el tamaño de las personas evaluadoras.} Las guías del área recomiendan cinco
personas evaluadoras. El equipo declara tres como mínimo viable y reportará el número efectivo. Si
solo se alcanzan tres, se declarará como amenaza a la validez de conclusión.

% =====================================================================
\section{Procedimiento}

\begin{enumerate}[leftmargin=1.8em]
  \item \textbf{Preparación del material fuente.} Se selecciona la transcripción completa de la
        entrevista \id{ENTR-04} y se verifica su anonimización: sin nombres propios, sin cargos que
        identifiquen de forma unívoca y sin referencias reconocibles a la organización.

  \item \textbf{Generación del conjunto A.} Se ejecuta el modelo con la consigna exacta:

        \begin{quote}
        \emph{``A partir del siguiente material fuente, redacta requisitos funcionales del sistema
        descrito, con los ocho atributos de la plantilla del sílabo.''}
        \end{quote}

        Se registran en \id{prompts\_llm/}: consigna literal, modelo y versión exacta, temperatura,
        top-p, top-k, semilla, fecha y hora de la consulta, y respuesta completa sin edición.

  \item \textbf{Preparación del conjunto B.} Se extraen del ERS los requisitos trazados a
        \id{EV-04}, presentados en el mismo formato que el conjunto A para impedir su
        identificación por la forma.

  \item \textbf{Normalización de formato.} Ambos conjuntos se transcriben a una plantilla común:
        misma tipografía, misma estructura de campos, identificadores neutros del tipo
        \id{R-001} a \id{R-0NN}. Se elimina toda marca que revele la procedencia.

  \item \textbf{Aleatorización y ciego.} Los requisitos de ambos conjuntos se mezclan en orden
        aleatorio con semilla registrada. La correspondencia entre identificador neutro y
        procedencia se conserva en un archivo separado que las personas evaluadoras no reciben.

  \item \textbf{Evaluación.} Cada persona evaluadora puntúa cada requisito de 1 a 5 en las cinco
        dimensiones, de forma independiente y sin comunicación con las demás.

  \item \textbf{Análisis.} Se aplica el plan de la Sección 8.

  \item \textbf{Desciego y discusión.} Se revela la procedencia y se analizan los desacuerdos entre
        personas evaluadoras y los casos extremos.
\end{enumerate}

% =====================================================================
\section{Instrumentos}

\begin{longtable}{|C{2.6cm}|L{5.0cm}|L{6.8cm}|}
\hline
\rowcolor{grisclaro}\textbf{Instrumento} & \textbf{Descripción} & \textbf{Ubicación} \\
\hline
\endhead
Material fuente & Transcripción anonimizada de \id{ENTR-04} & \id{02\_Evidencias/Transcripciones/} \\ \hline
Consigna & Texto literal, parámetros del modelo y respuesta & \id{06\_Experimento/prompts\_llm/} \\ \hline
Rúbrica & Cinco dimensiones, escala 1--5, con descriptor por nivel & \id{06\_Experimento/instrumentos/} \\ \hline
Hoja de puntuación & Una por persona evaluadora, en formato tabular & \id{06\_Experimento/instrumentos/} \\ \hline
Consentimiento & Formulario para las personas evaluadoras & \id{06\_Experimento/instrumentos/} \\ \hline
\caption{Instrumentos del estudio}
\end{longtable}

\noindent
\textbf{Descriptores de la rúbrica.} Cada dimensión se puntúa según la escala siguiente, común a las
cinco: 1, la propiedad está ausente; 2, presente de forma marginal; 3, presente con deficiencias
apreciables; 4, presente con deficiencias menores; 5, plenamente presente. Cada nivel se acompaña de
un ejemplo positivo y uno negativo en el documento de la rúbrica, con el fin de reducir la
variabilidad de criterio entre personas evaluadoras.

% =====================================================================
\section{Plan de análisis estadístico}

El plan se fija \emph{antes} de disponer de datos. No se modificará en función de los resultados
obtenidos.

\begin{enumerate}[leftmargin=1.8em]
  \item \textbf{Estadística descriptiva.} Media, mediana, desviación típica y rango intercuartílico
        por conjunto y por dimensión.

  \item \textbf{Prueba de normalidad.} Shapiro--Wilk sobre las diferencias apareadas de cada
        dimensión, con $\alpha = 0{,}05$.

  \item \textbf{Prueba de hipótesis.} Si las diferencias siguen distribución normal, prueba $t$ para
        muestras apareadas. En caso contrario, prueba de los rangos con signo de Wilcoxon.
        Contraste bilateral, $\alpha = 0{,}05$.

  \item \textbf{Corrección por comparaciones múltiples.} Se contrastan cinco dimensiones sobre el
        mismo conjunto de datos. Se aplica corrección de Holm--Bonferroni sobre los cinco valores
        $p$ obtenidos. Se reportan tanto los valores crudos como los corregidos.

  \item \textbf{Tamaño del efecto.} Cohen $d$ para muestras apareadas cuando se aplique la prueba
        paramétrica; Cliff $\delta$ cuando se aplique la no paramétrica. Se reporta con su intervalo
        de confianza al 95\%.

  \item \textbf{Acuerdo entre personas evaluadoras.} Coeficiente $\kappa$ de Cohen para cada par de
        personas evaluadoras y $\kappa$ de Fleiss para el conjunto completo. Se interpreta según los
        umbrales convencionales del área.

  \item \textbf{Potencia estadística.} Cálculo previo fijando $\alpha = 0{,}05$, potencia
        $1-\beta = 0{,}80$ y tamaño de efecto medio $d = 0{,}5$. Con esos parámetros, el tamaño
        muestral requerido para una prueba $t$ apareada bilateral es de aproximadamente 34 pares.
        \textbf{El diseño contempla 25 pares}, lo que sitúa la potencia alcanzable por debajo del
        objetivo. Esta limitación se declara de forma expresa y se reporta como amenaza a la validez
        de conclusión; no se incrementará artificialmente el tamaño muestral duplicando requisitos.
\end{enumerate}

\noindent
Todas las tablas y figuras del informe de resultados se producirán mediante los scripts versionados
en \id{06\_Experimento/scripts\_analisis/}. No se aceptarán cifras calculadas manualmente.

% =====================================================================
\section{Amenazas a la validez previstas}

\begin{longtable}{|C{2.6cm}|L{5.6cm}|L{6.2cm}|}
\hline
\rowcolor{grisclaro}\textbf{Categoría} & \textbf{Amenaza} & \textbf{Mitigación prevista} \\
\hline
\endhead
Constructo & La rúbrica de cinco dimensiones puede no capturar la calidad de un requisito en su totalidad & Las dimensiones proceden de la literatura del área y de los criterios transversales de la propia guía; se declara la limitación \\ \hline
Interna & El conjunto humano fue elaborado por el mismo equipo que ejecuta el estudio & Evaluación a ciegas, normalización de formato e identificadores neutros \\ \hline
Interna & Efecto de posición en el orden de presentación & Aleatorización con semilla registrada \\ \hline
Externa & Un único material fuente, un único dominio y un único modelo de lenguaje & Se declara. No se generalizará más allá del caso estudiado \\ \hline
Externa & Las personas evaluadoras son estudiantes, no profesionales en ejercicio & Se declara el perfil real; no se atribuirá al juicio la condición de experto profesional \\ \hline
Conclusión & Potencia estadística por debajo de $0{,}80$ con 25 pares & Se declara antes de ejecutar; se reporta el intervalo de confianza del tamaño del efecto y no solo el valor $p$ \\ \hline
Conclusión & Acuerdo bajo entre personas evaluadoras & Se reporta $\kappa$ con independencia de su valor; si resulta bajo, se discute como hallazgo y no se descarta \\ \hline
\caption{Amenazas a la validez identificadas antes de la ejecución}
\end{longtable}

% =====================================================================
\section{Consideraciones éticas}

Las personas evaluadoras firman un consentimiento específico que declara el propósito del estudio,
el carácter voluntario de su participación, la ausencia de compensación y el uso académico de sus
puntuaciones de forma agregada y anónima.

El material fuente sometido al modelo de lenguaje está anonimizado: no contiene nombres propios,
cargos identificativos ni referencias reconocibles a la organización. La persona participante
\id{ENTR-04} autorizó de forma expresa el uso de sus datos anonimizados en publicaciones
científicas revisadas por pares, conforme a la adenda ética de la segunda ronda.

Ninguna puntuación individual se publicará de forma atribuible a una persona evaluadora concreta.

% =====================================================================
\section{Compromiso de reporte íntegro}

El equipo se compromete de forma expresa a:

\begin{enumerate}[nosep,leftmargin=1.8em]
  \item Reportar el resultado obtenido con independencia de si confirma o refuta la hipótesis nula.
  \item No modificar el plan de análisis en función de los datos observados.
  \item Publicar los datos crudos de puntuación junto con los scripts que producen cada tabla.
  \item Declarar el número efectivo de personas evaluadoras y de requisitos evaluados.
  \item No redactar las secciones de resultados ni de discusión antes de disponer de datos
        analizados.
\end{enumerate}

\noindent
\fbox{\parbox{0.97\textwidth}{%
\textbf{Estado a la fecha de este documento.} El protocolo se encuentra \textbf{diseñado y
registrado}; el experimento \textbf{no ha sido ejecutado}. En consecuencia, el ERS/SRS de la Entrega
3 (2A) no contiene sección de resultados ni de discusión de este estudio. Ambas corresponden a la
Entrega 4 (2B), una vez recogidos y analizados los datos primarios.}}

\vspace{1.4cm}
\noindent
\begin{tabular}{p{7.0cm}p{7.0cm}}
\rule{6.4cm}{0.4pt} & \rule{6.4cm}{0.4pt} \\
Villafuerte Rosero Allan Noe & Ing. Gleiston Cicerón Guerrero Ulloa, PhD \\
Analista líder --- Equipo AHMRV & Docente responsable --- ISR-401 \\
\end{tabular}

\vspace{1.0cm}
\noindent Quevedo, 3 de agosto de 2026.

\end{document}
